% !TEX root = ../../ctfp-print.tex

\lettrine[lhang=0.17]{我}{们之前见过}, 当我们在范畴 $\cat{C}$ 中固定一个物件 $a$ 时, 映射
$\cat{C}(a, -)$ 是一个从 $\cat{C}$ 到 $\Set$ 的 (协变) 函子.
\[x \to \cat{C}(a, x)\]
(陪域是 $\Set$, 因为 hom-集 $\cat{C}(a, x)$ 是一个\emph{集合}.) 我们把这个映射叫作 hom-函子
--- 之前也定义过它对态射的作用.

现在让我们在这个映射里变动 $a$. 我们得到一个新的映射, 它把 hom-\emph{函子}
$\cat{C}(a, -)$ 赋给任意一个 $a$.
\[a \to \cat{C}(a, -)\]
它是从范畴 $\cat{C}$ 的物件到函子的映射, 而函子是函子范畴里的\emph{物件} (见
\hyperref[natural-transformations]{自然变换}中讲函子范畴的那一节). 我们用记号
$[\cat{C}, \Set]$ 表示从 $\cat{C}$ 到 $\Set$ 的函子范畴. 你可能还记得, hom-函子是最典型的
\hyperref[representable-functors]{可表函子}.

每当我们在两个范畴之间有一个物件到物件的映射, 都很自然地想问这样的映射是不是也是个函子.
换句话说, 我们能不能把一个范畴里的态射提升到另一个范畴里的态射. $\cat{C}$ 中的一个态射不过是
$\cat{C}(a, b)$ 的一个元素, 但函子范畴 $[\cat{C}, \Set]$ 中的态射是一个自然变换. 所以我们要
找的, 是一个把态射映射到自然变换的映射.

我们来找找看, 能不能为态射 $f \Colon a \to b$ 找到对应的自然变换. 先看看 $a$ 和 $b$ 被映到了
什么. 它们被映到两个函子: $\cat{C}(a, -)$ 和 $\cat{C}(b, -)$. 我们需要这两个函子之间的一个
自然变换.

诀窍在这里: 我们用米田引理:
\[[\cat{C}, \Set](\cat{C}(a, -), F) \cong F a\]
把泛型的 $F$ 换成 hom-函子 $\cat{C}(b, -)$. 于是得到:
\[[\cat{C}, \Set](\cat{C}(a, -), \cat{C}(b, -)) \cong \cat{C}(b, a)\]
\begin{figure}[H]
  \centering
  \includegraphics[width=0.6\textwidth]{images/yoneda-embedding.jpg}
\end{figure}

\noindent
这正是我们要找的、两个 hom-函子之间的自然变换, 只是拐了个小弯: 我们得到的对应, 一边是自然变换,
另一边是一个态射 --- $\cat{C}(b, a)$ 的一个元素 --- 而它的方向是 ``反'' 的. 不过这没关系;
它只说明我们正在考察的这个函子是逆变的.

\begin{figure}[H]
  \centering
  \includegraphics[width=0.65\textwidth]{images/yoneda-embedding-2.jpg}
\end{figure}

\noindent
实际上, 我们赚到的比想要的还多. 从 $\cat{C}$ 到 $[\cat{C}, \Set]$ 的这个映射不仅是一个逆变
函子 --- 它还是一个\emph{全忠实} (fully faithful) 函子. 全与忠实是函子的性质, 描述它们如何映射
hom-集.

一个\emph{忠实} (faithful) 函子在 hom-集上是 \newterm{injective 单射} 的, 也就是说它把不同的
态射映到不同的态射. 换句话说, 它不会把它们合并.

一个\emph{全} (full) 函子在 hom-集上是 \newterm{surjective 满射} 的, 也就是说它把一个 hom-集
\emph{映上} (onto) 另一个 hom-集, 完全覆盖后者.

一个全忠实的函子 $F$ 在 hom-集上是一个 \newterm{bijection 双射} --- 两个集合的所有元素一一
配对. 对源范畴 $\cat{C}$ 中每一对物件 $a$ 和 $b$, $\cat{C}(a, b)$ 与 $\cat{D}(F a, F b)$
之间存在一个双射, 其中 $\cat{D}$ 是 $F$ 的目标范畴 (在我们的情形里就是函子范畴,
$[\cat{C}, \Set]$). 注意, 这并不意味着 $F$ 在\emph{物件}上是双射. $\cat{D}$ 中可能有些物件
不在 $F$ 的像里, 关于这些物件的 hom-集我们什么也说不出来.

\section{嵌入}

我们刚刚描述的 (逆变) 函子 --- 把 $\cat{C}$ 中的物件映到 $[\cat{C}, \Set]$ 中的函子的那个
函子:
\[a \to \cat{C}(a, -)\]
定义了 \newterm{Yoneda embedding 米田嵌入}. 它把范畴 $\cat{C}$ \emph{嵌入}到了函子范畴
$[\cat{C}, \Set]$ 内部 (严格来说, 由于逆变, 被嵌入的是范畴 $\cat{C}^{op}$). 它不仅把
$\cat{C}$ 中的物件映成函子, 还忠实地保留了它们之间的所有连接.

这是一个非常有用的结论, 因为数学家们对函子的范畴了解很多, 尤其是那些陪域为 $\Set$ 的函子.
把一个任意的范畴 $\cat{C}$ 嵌入函子范畴, 能让我们对它获得大量洞见.

当然, 米田嵌入也有对偶版本, 有时叫作余米田嵌入. 注意我们本可以从固定每个 hom-集的目标物件
(而不是源物件) 开始, 即 $\cat{C}(-, a)$. 那会给出一个逆变的 hom-函子. 从 $\cat{C}$ 到
$\Set$ 的逆变函子就是我们熟悉的预层 (参见, 比如说,
\hyperref[limits-and-colimits]{极限与余极限}). 余米田嵌入定义了一个范畴 $\cat{C}$ 在预层范畴
中的嵌入. 它对态射的作用由下式给出:
\[[\cat{C}, \Set](\cat{C}(-, a), \cat{C}(-, b)) \cong \cat{C}(a, b)\]
同样, 数学家们对预层的范畴了解很多, 所以能把任意范畴嵌入其中, 是一桩大好事.

\section{在 Haskell 中的应用}

在 Haskell 里, 米田嵌入可以表示成一个同构: 一边是 reader 函子之间的自然变换, 另一边是函数
(而且方向相反):

\begin{snipv}
forall x. (a -> x) -> (b -> x) \ensuremath{\cong} b -> a
\end{snipv}
(记住, reader 函子等价于 \code{((->) a)}.)

这个等式的左边是一个多态函数: 给定一个从 \code{a} 到 \code{x} 的函数, 再加上一个类型为
\code{b} 的值, 它就能产出一个类型为 \code{x} 的值 (我这里是在做未柯里化 --- 把 \code{b -> x}
外面的那层括号去掉.) 要对所有 \code{x} 都做到这一点, 唯一的办法就是我们的函数知道如何把一个
\code{b} 转换成 \code{a}. 它必定偷偷地掌握着一个 \code{b -> a} 的函数.

给定这样一个转换器 \code{btoa}, 我们就可以定义左边 --- 叫它 \code{fromY} --- 为:

\src{snippet01}
反过来, 给定函数 \code{fromY}, 只要用恒等函数调用它, 就能把转换器找回来:

\src{snippet02}
这就在 \code{fromY} 类型的函数与 \code{btoa} 之间建立了一个双射.

看待这个同构的另一种方式是: 它是一个从 \code{b} 到 \code{a} 的函数的 \acronym{CPS} 编码.
参数 \code{a -> x} 是一个续体 (即处理器). 结果是一个从 \code{b} 到 \code{x} 的函数, 当它带着
一个类型为 \code{b} 的值被调用时, 就会执行那个与被编码函数前组合过的续体.

米田嵌入还解释了 Haskell 中数据结构的某些替代表示. 特别地, 它为 \code{Control.Lens} 库里的
透镜 (lens) 提供了一个
\urlref{https://bartoszmilewski.com/2015-07-13/from-lenses-to-yoneda-embedding/}
{非常好用的表示}.

\section{预序的例子}

这个例子是 Robert Harper 提议的. 它是米田嵌入在一个由预序定义的范畴上的应用. 预序是一个集合,
它的元素之间有一个序关系, 传统上写作 $\leqslant$ (小于或等于). 预序 (preorder) 里的这个 ``预''
字之所以在, 是因为我们只要求这个关系是传递的和自反的, 而不一定反对称 (所以其中可能出现环).

带有预序关系的集合导出一个范畴. 物件就是这个集合的元素. 从物件 $a$ 到 $b$ 的态射, 当两个物件
无法比较、或者 $a \leqslant b$ 不成立时, 它不存在; 当 $a \leqslant b$ 成立时, 它存在, 并且从
$a$ 指向 $b$. 从一个物件到另一个物件永远不会有多于一个态射. 因此这种范畴里的任何 hom-集要么是
空集, 要么是单例集合. 这样的范畴叫作\emph{薄}范畴.

很容易说服自己这个构造确实是一个范畴: 箭头是可组合的, 因为若 $a \leqslant b$ 且 $b \leqslant c$, 则
$a \leqslant c$; 而且组合是结合的. 我们也有恒等箭头, 因为每个元素都 (小于或)
等于自身 (这是底层关系的自反性).

现在我们可以把余米田嵌入应用于预序范畴. 特别地, 我们关心它对态射的作用:
\[[\cat{C}, \Set](\cat{C}(-, a), \cat{C}(-, b)) \cong \cat{C}(a, b)\]
右边的 hom-集非空, 当且仅当 $a \leqslant b$ --- 此时它是一个单例集合. 于是, 若
$a \leqslant b$, 左边就存在唯一一个自然变换. 否则就没有自然变换.

那么, 预序中 hom-函子之间的自然变换是什么? 它应该是在集合 $\cat{C}(-, a)$ 与 $\cat{C}(-, b)$
之间的一族函数. 在预序里, 这些集合要么是空的, 要么是单例的. 我们来看看手上都有哪些函数可用.

从空集到自身有一个函数 (作用在空集上的恒等函数), 从空集到单例集合有一个 \code{absurd} 函数
(它什么都不做, 因为它只需要对空集的元素有定义, 而这样的元素一个也没有), 从单例集合到自身也有一个
函数 (作用在单元素集合上的恒等函数). 唯一被禁止的组合是从单例集合映到空集 (这样的函数作用在那
唯一的元素上, 取值该是什么呢?).

所以我们的自然变换绝不会把单例 hom-集连到空 hom-集. 换句话说, 若 $x \leqslant a$ (单例 hom-集
$\cat{C}(x, a)$), 那么 $\cat{C}(x, b)$ 就不可能是空的. $\cat{C}(x, b)$ 非空意味着 $x$ 小于或
等于 $b$. 于是这个自然变换的存在, 要求对每个 $x$, 若 $x \leqslant a$ 则 $x \leqslant b$.
\[\text{for all } x, x \leqslant a \Rightarrow x \leqslant b\]
另一方面, 余米田告诉我们, 这个自然变换的存在等价于 $\cat{C}(a, b)$ 非空, 亦即等价于
$a \leqslant b$. 合在一起, 我们得到:
\[a \leqslant b \text{ if and only if for all } x, x \leqslant a \Rightarrow x \leqslant b\]
我们本来也能直接得出这个结果. 直觉是, 若 $a \leqslant b$, 那么所有低于 $a$ 的元素必然也低于
$b$. 反过来, 在右边把 $x$ 换成 $a$, 就推出 $a \leqslant b$. 但你得承认, 通过米田嵌入走到这个
结果, 要刺激得多.

\section{自然性}

米田引理在自然变换的集合与 $\Set$ 中的一个物件之间建立了同构. 自然变换是函子范畴
$[\cat{C}, \Set]$ 中的态射. 任意两个函子之间的自然变换组成的集合, 是这个范畴里的一个 hom-集.
米田引理就是这个同构:
\[[\cat{C}, \Set](\cat{C}(a, -), F) \cong F a\]
这个同构对 $F$ 和 $a$ 两者都是自然的. 换句话说, 它对 $(F, a)$ 是自然的, 而这个配对取自积范畴
$[\cat{C}, \Set] \times \cat{C}$. 注意, 我们现在把 $F$ 当作函子范畴里的一个 \newterm{object
物件}.

我们想一想这意味着什么. 自然同构是两个函子之间一个可逆的\emph{自然变换}. 而我们这个同构的右边
确实是一个函子: 它是从 $[\cat{C}, \Set] \times \cat{C}$ 到 $\Set$ 的函子. 它作用在配对
$(F, a)$ 上的结果是一个集合 --- 即把函子 $F$ 在物件 $a$ 处求值的结果. 这叫作求值函子 (evaluation
functor).

左边也是一个函子, 它把 $(F, a)$ 送到自然变换的集合
$[\cat{C}, \Set](\cat{C}(a, -), F)$.

要证明它们真的是函子, 我们还得定义它们对态射的作用. 可是从配对 $(F, a)$ 到 $(G, b)$ 的态射
是什么? 它是一个态射的配对 $(\Phi, f)$; 第一个是函子之间的态射 --- 一个自然变换 --- 第二个是
$\cat{C}$ 中的普通态射.

求值函子接受这个配对 $(\Phi, f)$, 把它映成两个集合 $F a$ 和 $G b$ 之间的一个函数. 我们很容易
用 $\Phi$ 在 $a$ 处的分量 (它把 $F a$ 映到 $G a$) 以及被 $G$ 提升后的态射 $f$, 构造出这样一个
函数:
\[(G f) \circ \Phi_a\]
注意, 由于 $\Phi$ 的自然性, 这和下面这个是同一个函数:
\[\Phi_b \circ (F f)\]
我不打算证明整个同构的自然性 --- 一旦界定了那些函子是什么, 证明就相当机械. 它从这个事实推出:
我们的同构是由函子和自然变换搭建起来的. 它根本没有出错的可能.

\section{挑战}

\begin{enumerate}
  \tightlist
  \item
        用 Haskell 表达余米田嵌入.
  \item
        证明我们在 \code{fromY} 与 \code{btoa} 之间建立的双射是一个同构 (两个映射互为逆).
  \item
        对幺半群算出米田嵌入. 幺半群唯一的那个物件对应于哪个函子? 幺半群的态射又对应于哪些
        自然变换?
  \item
        \emph{协变}的米田嵌入应用于预序会得到什么? (此题由 Gershom Bazerman 提议.)
  \item
        米田嵌入可以用来把任意函子范畴 $[\cat{C}, \cat{D}]$ 嵌入函子范畴
        $[[\cat{C}, \cat{D}], \Set]$. 想想它对态射是如何起作用的 (在这里, 态射就是自然变换).
\end{enumerate}
